\documentclass[10pt,a4paper]{amsart}
\usepackage[margin=19mm]{geometry}
\usepackage{amsmath,amssymb,mathtools}
\usepackage[scr=rsfs,cal=euler]{mathalfa}
\usepackage{xfrac}
\usepackage[unicode,hidelinks,bookmarks=false]{hyperref}
\newcommand{\ie}{i.e.\ }
\newcommand{\cf}{cf.\ }
\newcommand{\resp}{respectively\ }
\newcommand{\Subsection}{\subsection}
\providecommand{\nobreakdash}{\nobreak\hbox{-}}
\setlength{\emergencystretch}{2em}
\multlinegap0pt
\usepackage{amsrefs}
\usepackage{framed}
\usepackage{array}
\numberwithin{equation}{section}
\newtheorem{theorem}{Theorem}
\newtheorem{lemma}{Lemma}
\newtheorem{conj}{Conjecture}
\newtheorem{coro}{Corollary}
\newcommand{\eps}{\varepsilon}
\newcommand{\Z}{\mathbb Z}
\newcommand{\N}{\mathbb N}
\newcommand{\R}{\mathbb R}
\newcommand{\C}{\mathbb C}
\newcommand{\Q}{\mathbb Q}
\def\legendre#1#2{\genfrac(){0.2pt}0{#1}{#2}}
\newcommand{\minilegendre}[2]{\genfrac{(}{)}{}{}{#1}{#2}}
\begin{document}
\title[On series involving binom4kk^-1]{On series involving $\binom{4k}{k}^{-1}$}
\author{Roman Le Lan}
\date{}
\maketitle
\noindent\emph{Source and licence.} On series involving $\binom{4k}{k}^{-1}$. Version arXiv:2609.22304v1 (CC BY 4.0). This material is distributed under the Creative Commons Attribution 4.0 International licence, http://creativecommons.org/licenses/by/4.0/, as recorded in the arXiv OAI licence metadata for this exact version and shown on the abstract page https://arxiv.org/abs/2609.22304v1. Full rights notice: licenses/arxiv-2609.22304-CC-BY-4.0.txt.\par\smallskip
\noindent\textit{Editorial scope.} Counted unit: the original \texttt{lemma} environment \texttt{lem3} at source lines 113--119 together with its complete proof at lines 120--152; the delivered document keeps source lines 75--153 so that every referenced label and every citation resolves inside it. Native editable source is the paper's own concatenated TeX; the AMS wrapper is editorial formatting only. No publisher class, upstream script or unrelated artwork is redistributed.
\section{Preliminary results}
The proofs of the theorems in this note rely on a few intermediate results. The first lemma provides a decomposition for quadratic polynomials.
\begin{lemma}
    For every real polynomial $P$ of degree at most $2$, there exists a triple $(\alpha,\beta,\gamma)\in\R^3$ such that:
    \[P(X)=\alpha (3X-1)(3X-2)+\beta X(X+1)+\gamma X(4X-1)\]
\end{lemma}
Note that $\alpha$, $\beta$, and $\gamma$ can be readily determined by solving a system of linear equations. We will compute these values as needed. 
\begin{proof}
    Consider the vector space $\R_2[X]$ of polynomials of degree at most $2$. Define the polynomials:
    \[P_1=(3X-1)(3X-2),\quad P_2=X(X+1),\quad P_3=X(4X-1).\]
    It is straightforward to verify that $(a,b,c)=(0,0,0)$ is the unique solution to the equation:
    \[aP_1+bP_2+cP_3=0\]
    Therefore, $P_1$, $P_2$, and $P_3$ are linearly independent in $\R_2[X]$.
    
    Moreover, since $\dim\R_2[X]=3$, the set $(P_1,P_2,P_3)$ forms a basis for $\R_2[X]$. 

    Consequently, every polynomial $Q\in\R_2[X]$ can be uniquely written as a linear combination of $P_1$, $P_2$, and $P_3$.
\end{proof}
Our second lemma introduces a useful transformation used in the proof of Lemma 2.
\begin{lemma}
    For all integers $n,m\ge1$, 
    \begin{equation}
     \frac{H_{n+m-1}-H_{n-1}}{n\binom{n+m-1}{n}}=-\frac{\partial}{\partial x}B(x,y)\Bigg|_{\substack{x=n\\y=m}}.
    \label{lem2}
    \end{equation}
\end{lemma}
\begin{proof}
    Observe that
    \begin{align*}
        \frac{\partial}{\partial x}B(x,y)&=\frac{\partial}{\partial x}\frac{\Gamma(x)\Gamma(y)}{\Gamma(x+y)}\\&=\frac{\Gamma(x)\Gamma(y)}{\Gamma(x+y)}(\psi(x)-\psi(x+y))\\&=B(x,y)(\psi(x)-\psi(x+y))
    \end{align*}
    where $\psi(x)=\frac{d}{dx}\log{\Gamma(x)}$. Evaluating at $(x,y)=(n,m)$ for integers $n,m\ge1$, we obtain:
    \begin{align*}
        \frac{\partial}{\partial x}B(x,y)\Bigg|_{\substack{x=n\\y=m}}&=B(n,m)(\psi(n)-\psi(n+m))\\&=\frac{1}{n\binom{n+m-1}{n}}(H_{n-1}-H_{n+m-1})
    \end{align*}
    since $\psi(n)=H_{n-1}-\gamma$, where $\gamma$ is the Euler-Mascheroni constant. This proves the lemma.
\end{proof}
The third lemma will be useful at the end of the proof of Theorem 2.

\begin{lemma}
    We have the following identity:
    \begin{equation}
    \int_0^1\frac{\log(1-t)}{3t^2-6t+4}\,dt=-\frac{\pi\sqrt{3}}{18}\log3-\frac{5}{8}K.
    \label{lem3}
    \end{equation}
\end{lemma}

\begin{proof}
    Let 
    \[I=\int_0^1\frac{\log(1-t)}{3t^2-6t+4}\,dt.\]
    Make the change of variable $t=1-x$. Then,
    \[I=\int_0^1\frac{\log(x)}{3x^2+1}\,dx.\]
    Now, making the change of variable $x\sqrt3=\tan\theta$, we obtain:
    \begin{align}
        I&=\frac{\sqrt3}{3}\int_0^{\pi/3}\log\left(\frac{\tan\theta}{\sqrt3}\right)\,d\theta\nonumber\\
        &=-\frac{\pi\sqrt{3}}{18}\log3+\frac{\sqrt3}{3}\left(\int_0^{\pi/3}\log\left(\sin\theta\right)\,d\theta-\int_0^{\pi/3}\log\left(\cos\theta\right)\,d\theta\right)\nonumber\\
        &:=-\frac{\pi\sqrt{3}}{18}\log3+\frac{\sqrt3}3(J_1-J_2)
        \label{I=J1-J2}
    \end{align}
    where $J_1$ and $J_2$ are the respective integrals in the second line. Recall the integral definition of the Clausen function:
    \[\text{Cl}_2(z)=-\int_0^z\log\left|2\sin\frac{\theta}{2}\right|\,d\theta.\]
    Making the change of variable $\theta=u/2$ in the expression for $J_1$, it is straightforward to verify that:
    \begin{equation}
        J_1=-\frac12\text{Cl}_2(2\pi/3)-\frac\pi3\log2.
    \label{J1}
    \end{equation}
    Moreover, making the change of variable $u=\pi/2-\theta$ in the expression for $J_2$, we find:
    \[J_2=\int_{\pi/6}^{\pi/2}\log{(\sin u)}\,du.\]
    This can be evaluated in the same manner as $J_1$. Thus, after some simplification,
    \begin{equation}
    J_2=\frac12\text{Cl}_2(\pi/3)-\frac\pi3\log2.
    \label{J2}
    \end{equation}
    Since \[\text{Cl}_2(\pi/3)=\frac32\text{Cl}_2(2\pi/3)=\frac{3\sqrt3}{4}K,\]
    substituting \eqref{J1} and \eqref{J2} into \eqref{I=J1-J2}, we obtain:
    \begin{align*}
        I&=-\frac{\pi\sqrt{3}}{18}\log3-\frac{\sqrt3}3\cdot\frac{5\sqrt3}{8}K\\&=-\frac{\pi\sqrt{3}}{18}\log3-\frac{5}{8}K.
    \end{align*}
    This proves Lemma 3.
\end{proof}

\end{document}
